\documentclass[handout]{beamer}

\usetheme{Warsaw}

\usepackage[brazil]{babel}
\usepackage[latin1]{inputenc}
\usepackage{times}
\usepackage[T1]{fontenc}
\usepackage{epsfig}
\usepackage{listings}
\usepackage{amsfonts}
\usepackage{multirow}
\usepackage{multicol}
\usepackage{tikz}
\usepackage{textpos}
\usepackage{framed, multicol}
\usepackage{wrapfig}
\usepackage{array}
\usepackage{booktabs}
\usepackage{verbatim}
\usepackage{tikz}
\usepackage{tikz-uml}


%\usepackage{dcolumn}
%\newcolumntype{.}{D{.}{.}{-1}}
%\newcolumntype{d}[1]{D{.}{.}{#1}}
%\usetheme{Berkeley}

\newcommand{\be}{\begin{enumerate}[<+->]}
\newcommand{\ee}{\end{enumerate}}

\newcommand{\bq}{\begin{quote}}
\newcommand{\eq}{\end{quote}}

\newcommand{\bd}{\begin{description}[<+->]}
\newcommand{\ed}{\end{description}}

\newcommand{\bi}{\begin{itemize}[<+->]}
\newcommand{\ei}{\end{itemize}}

\newcommand{\bbl}{\begin{block}{Exemplo: Tipos primitivos do Pascal}
        \scriptsize
        \centering
        }
\newcommand{\ebl}{\end{block}}




%%%%%%%%%%%%%%%%%%%%%%%%%%%%%%%%%%%%%%%%%%%%%%%%%%%%%%%%%%%%%%%%
% Configuração do Tiks
%%%%%%%%%%%%%%%%%%%%%%%%%%%%%%%%%%%%%%%%%%%%%%%%%%%%%%%%%%%%%%%%
\usetikzlibrary{arrows}

\tikzstyle{bloco}=[draw, minimum size=2em] %fill=blue!20,
\tikzstyle{linha} = [pin edge={to-,thin,black}]



\title[Programação Modular]
{%
    \textit{Erros de Software e TDD}%
}
\author[Prof. Hugo de Paula]
{
    Prof.~Hugo~de~Paula
}
\institute[DCC / PUC Minas]
{\epsfig{file=puclogo_small_bw,width=1.5cm} \\
  \textsc{Pontifícia Universidade Católica de Minas Gerais}\\
    Departamento de Ciência da Computação
}
\date[]{}

\lstset{language=Java,
   basicstyle=\scriptsize,
   commentstyle=\color{red},
   showstringspaces=false,
   numbers=none,
   numberstyle=\tiny}

\begin{document}


\selectlanguage{brazil}

\begin{frame}
   \titlepage
\end{frame}

%\addtobeamertemplate{frametitle}{}{%
%    \begin{tikzpicture}[remember picture,overlay]
%    \node[anchor=north east,yshift=2pt] at (current page.north east) {\epsfig{file=puclogo_small_bw,width=1.2cm}};
%    \end{tikzpicture}}


%\addtobeamertemplate{frametitle}{}{%
    %\begin{tikzpicture}[node distance=0cm, remember picture, overlay, every node/.style={inner sep=0,outer sep=0, node distance=0cm, baseline=0cm}]
    %\node[anchor=north east] at (current page.north east) {\epsfig{file=puclogo_small_bw,width=1cm}};
    %\end{tikzpicture}}


%\logo{\includegraphics[height=0.8cm]{puclogo_small_bw.pdf}\vspace{220pt}}


\begin{frame}
   \frametitle{Sumário}
   \tableofcontents[pausesections]
\end{frame}

%\AtBeginSection[] % Do nothing for \section*
%{
%\begin{frame}<beamer>
%\frametitle{Outline}
%\tableofcontents[currentsection]
%\end{frame}}


\addtobeamertemplate{frametitle}{}{%
   \begin{textblock*}{10mm}(.9945\textwidth,-2.01cm)
    \includegraphics[height=1cm]{puclogo_small_bw.pdf}
   \end{textblock*}
}

\section{Princípios do teste de software}

\begin{frame}[fragile]{Índice de Massa Corporal}

Considere o exemplo a seguir que calcula o Índice de Massa Corporal (IMC) de uma pessoa, dados seu peso e altura.

\begin{lstlisting}
public class CalculadoraIMC {
  private double peso, altura, imc;
  
  public CalculadoraIMC(double peso, double altura) {
    this.peso = peso;
    this.altura = altura;
  }
  public void calcular() {
    this.imc = peso / (altura * altura);
  }
  public boolean comSobrepeso() {
    return (imc > 25);
  }
  public double getImc() {
    return this.imc;
  }
}
\end{lstlisting}

\end{frame}


\begin{frame}[fragile]{Índice de Massa Corporal}

\begin{itemize}
\item O estado do objeto corresponde ao \lstinline|peso|, à \lstinline|altura| e ao \lstinline|imc|.
\item O comportamento do objeto é ditado pelos métodos \lstinline|calcular()| e \lstinline|comSobrepeso()|.
\item O método \lstinline|getImc()| é apenas um método de acesso.
\end{itemize}

\begin{lstlisting}
public static void main(String[] args) {
  CalculadoraIMC calcImc = new CalculadoraIMC(70, 1.60);
  calcImc.calcular();
    
  System.out.println("O IMC e: " + calcImc.getImc());
 
  if (calcImc.comSobrepeso()) {
    System.out.println("Tem sobrepeso.");
  }
}
\end{lstlisting}

\end{frame}

\subsection{Tipos de erro de software} 

\begin{frame}{\textit{Debugging}}

\begin{block}{\textit{Debugging}}
Processo utilizado para remover erros de programação (\textit{bugs}).
\end{block}

Passos para o \textit{debugging}:

\begin{enumerate}
\item Detectar o erro.
\item Localizar o erro.
\item Solucionar o erro.
\end{enumerate}

Tipos de erros de programação:

\begin{itemize}
\item Erros de sintaxe / de compilação.
\item Erros em tempo de execução.
\item Erros semânticos ou de lógica.
\end{itemize}
\end{frame}

\begin{frame}[fragile]{Erros de sintaxe}

\begin{itemize}
\item Violam normas gramaticais da linguagem.
\item São capturados previamente pelos compiladores ou interpretadores.
\item São identificadas pela IDE.
\end{itemize}

\begin{lstlisting}
  public void calcular() {
    this.imc = peso / (altura * altura),
  }
\end{lstlisting}
\end{frame}


\begin{frame}[fragile]{Erros em tempo de execução}

\begin{itemize}
\item Ocorrem durante a execução do programa.
\item Não podem ser capturados pelo compilador ou pela IDE. 
\item Em algumas IDEs, o erro pode ser alertado. Exemplo: Variável não inicializada.
\item Devem ser antecipados e tratados.
\end{itemize}

\begin{lstlisting}
  public void calcular() {
    this.imc = peso / (altura * altura);
  }
\end{lstlisting}

Se a altura for igual a 0 (zero), ocorrerá o erro ``\textit{/ by zero}'' (divisão por zero).
\end{frame}

\begin{frame}[fragile]{Erros semânticos ou de lógica}

\begin{itemize}
\item Os mais difíceis de detectar. 
\item Programa irá ``funcionar'', mas a saída estará incorreta.
\end{itemize}

\begin{lstlisting}
  public void calcular() {
    this.imc = (peso * peso) / altura;
  }
\end{lstlisting}

\end{frame}


\subsection{Princípios do teste de software}

\begin{frame}{Princípios do teste de software}

\begin{itemize}
\item Objetivo do teste é verificação e validação.
\end{itemize}

\begin{block}{Verificação}
objetiva responder se o sistema foi construído corretamente.
\end{block}

\begin{block}{Validação}
tenta determinar se foi construído o sistema certo.
\end{block}

\end{frame}

\begin{frame}{Custo do teste de software}

\begin{itemize}
\item O custo do erro aumenta, quanto mais se demora a detectá-lo.
\end{itemize}

{\centering \includegraphics[width=.73\textwidth]{custoTeste.pdf}

{\tiny Adaptado de: \textit{Beginning Java Programming (2015) - Bart Baesens and Aimée Backiel}; Capítulo 1, pg. 7}}
\end{frame}



\section{Test-driven development (TDD)}

\subsection{TDD: o que é} 

\begin{frame}{\textit{Test-driven development}}

\begin{block}{\textit{Test-driven development (TDD)}}
Ou desenvolvimento dirigido por teste é uma metodologia de desenvolvimento de software que enfatiza o teste.
\end{block}

``Escreva o código de teste antes e então escreva o código operacional e depure o programa até que ele passe no teste.''
\begin{itemize}
\item TDD é uma metodologia de desenvolvimento, não de teste.

\item Software é desenvolvido em pequenas iterações.

\item Testes de unidade são desenvolvidos antes do código.
\end{itemize}
\end{frame}

\subsection{TDD: como funciona}

\begin{frame}{TDD: como funciona}

\begin{enumerate}
\item Adicione um teste.
\begin{itemize}
\item Estórias do usuário ajudam a entender os requisitos.
\end{itemize}
\item Execute todos os testes e veja o novo falhar.
\begin{itemize}
\item Garante que o mecanismo de teste está funcionando e que o teste não irá passar por engano. 
\end{itemize}
\item Escreva o código.
\begin{itemize}
\item Apenas o código que foi projetado para passar no teste.
\item Nenhuma funcionalidade adicional deve ser incluída, pois não seria testada.
\end{itemize}
\end{enumerate}

\end{frame}



\begin{frame}{TDD: como funciona}

\begin{enumerate}
\setcounter{enumi}{3}
\item Execute os testes e veja eles serem bem sucedidos.
\begin{itemize}
\item Se o teste passou, tem-se confiança de que o código atende aos requisitos testados. 
\end{itemize}
\item Refatore o código.
\begin{itemize}
\item Limpe o código e reexecute os testes para garantir que nada foi quebrado.
\end{itemize}
\end{enumerate}

Repita o processo.

\end{frame}


\section{Teste de unidade e JUnit}

\subsection{Teste de unidade}

\begin{frame}{Teoria do teste de unidade}

\begin{block}{Teste unitário (ou teste de unidade)}
Para cada trecho de código operacional -- uma classe ou um método, o código operacional deve estar pareado com um código de ``teste unitário''.
\end{block} 

\vspace{5mm}

\textit{``If you can't write a test for what you are about to code, then you shouldn't even be thinking about coding.''}\footnote{George, Boby; Williams, Laurie. \textit{A Structured experiment of test-driven development}. 2003.}
\end{frame}



\begin{frame}{Teste de unidade}

\begin{itemize}
\item O código do teste unitário chama o código operacional a partir da sua interface pública, de diversas formas e verifica os resultados.
\item O teste não precisa ser exaustivo -- código de teste já agrega um grande valor, mesmo nos casos centrais.
\item Testes unitários são uma forma padronizada de manter os testes em paralelo com o desenvolvimento do programa.
\item Testes de unidade são fáceis de executar.
\begin{itemize}
\item uma vez configurados, eles produzem \textit{feedback} rápido para você testar suas ideias.
\end{itemize}
\end{itemize}
\end{frame}

\begin{frame}{Tipos de teste de unidade}
\begin{itemize}
\item \textbf{Básicos}: entradas simples e óbvias que devem funcionar. São adicionados primeiro.

\item \textbf{Avançados}: casos de teste mais complexos. Normalmente feitos posteriormente quando um maior conhecimento do problema leva a identificar potenciais casos estranhos.

\item \textbf{Borda/fronteira}: são casos simples mas que representam condições extremas -- \lstinline|string| vazia, lista vazia, etc.

\end{itemize}

\end{frame}

\subsection{JUnit com Eclipse}

\begin{frame}{JUnit com Eclipse}
\begin{itemize}
\item JUnit é um sistema de teste de unidade em Java muito popular e bem integrado ao Eclipse.
\item Para criar um caso de teste:
\begin{itemize}
\item Clique com o botão direito na classe que deseja testar.
\item Selecione \texttt{New... JUnit Test Case}.
\item Se a classe se chama \lstinline|Shulambs|, será criada uma classe chamada \lstinline|ShulambsTest|.
\end{itemize}
\end{itemize}

\end{frame}

\begin{frame}{JUnit.jar -- adicionando ao projeto}
\begin{itemize}
\item JUnit depende de classes disponíveis no arquivo \texttt{junit.jar}, já instalado como plugin do Eclipse.

\item Ao criar uma classe de teste, o Eclipse irá perguntar se deseja adicionar o arquivo ao projeto automaticamente.

\item Caso tenha que adicionar o arquivo manualmente, \texttt{Project}, selecione \texttt{properties... Java Build Path: Libraries.} Então utilize o botão \texttt{Add Jar} para adicionar o arquivo, que está na pasta \texttt{plugins}.
\end{itemize}
\end{frame}


\subsection{Exemplo de caso de teste JUnit}

\begin{frame}{Exemplo JUnit: classe \texttt{Conjunto}}

A clase \lstinline|Conjunto| é uma coleção

\begin{itemize}
\item possui até $N$ objetos.
\item possui as operações básicas:
\end{itemize}

\begin{center}
\begin{tabular}{|l|l|}
\hline
\lstinline|adicionar| & adiciona um item \\ \hline
\lstinline|contem|    & verifica se item no conjunto \\ \hline
\lstinline|tamanho|   & número de itens \\ \hline
\end{tabular}
\end{center}

\end{frame}


\begin{frame}{Exemplo JUnit: lista de testes}

\begin{itemize}
\item Decidindo como testar:
\begin{itemize}
\item $tamanho=0$ se \lstinline|Conjunto| está vazio.
\item $tamanho=N$ após adicionar $N$ elementos.
\item adicionar um elemento que já existe no \lstinline|Conjunto| não aumenta seu tamanho.
\item lança uma exceção se \lstinline|Conjunto| já está lotado.
\end{itemize}
\item cada teste verifica uma característica específica de \lstinline|Conjunto|
\end{itemize}
\end{frame}


\begin{frame}[fragile]{Exemplo JUnit: primeiro teste}

\begin{tabular}{|c|c|}
\hline
\begin{minipage}{.35\textwidth}
\begin{lstlisting}
public class Conjunto {

}
\end{lstlisting}
\end{minipage} &
\begin{minipage}{.55\textwidth}
\begin{lstlisting}
public class ConjuntoTest {
   @Test 
   public void testConjuntoVazio() {
      Conjunto c = new Conjunto();
      assertEquals(0, c.tamanho());
   }
}
\end{lstlisting}
\end{minipage} \\ \hline
\end{tabular}

\begin{itemize}
\item Não compila porque \lstinline|tamanho()| não foi definido.
\end{itemize}

\end{frame}



\begin{frame}[fragile]{Exemplo JUnit: Testando o método \texttt{tamanho()}}

\textcolor[rgb]{.5,0,0}{\bf BARRA VERMELHA}: código falhou no teste.

\begin{tabular}{ll}
\begin{minipage}{.45\textwidth}
\begin{lstlisting}
public class Conjunto {
   public int tamanho() {
      return 42;
   }
}
\end{lstlisting} 
\end{minipage} & 
\includegraphics[width=.45\textwidth]{falhaConjunto.png} \\
\end{tabular}

\textcolor[rgb]{0,.5,0}{\bf BARRA VERDE}: código passou no teste.

\begin{tabular}{ll}
\begin{minipage}{.45\textwidth}
\begin{lstlisting}
public class Conjunto {
   public int tamanho() {
      return 0;
   }
}
\end{lstlisting} 
\end{minipage} & 
\includegraphics[width=.45\textwidth]{passaConjunto.png} \\
\end{tabular}

\end{frame}




\begin{frame}[fragile]{Exemplo JUnit: testando adicionar itens}

\begin{tabular}{|c|c|}
\hline
\begin{minipage}{.45\textwidth}
\begin{lstlisting}[basicstyle=\tiny]
public class Conjunto {
   public int tamanho() {
      return 0;
   }
   public void adicionar(Object o) {
      
   }
}
\end{lstlisting}
\end{minipage} &
\begin{minipage}{.5\textwidth}
\begin{lstlisting}[basicstyle=\tiny]
public class ConjuntoTest {
   @Test
   public void testConjuntoVazio() {
      Conjunto c = new Conjunto();
      assertEquals(0, c.tamanho());
   }

   @Test
   public void testAdicionarUm() {
      Conjunto c = new Conjunto();
      c.adicionar(new Object());
      assertEquals(1,c.tamanho());
   }
}
\end{lstlisting}
\end{minipage} \\ \hline
\end{tabular}

\vspace{2mm}

\textcolor[rgb]{.5,0,0}{\bf Falhou} como esperado, pois o tamanho não está sendo contado.

\end{frame}






\begin{frame}[fragile]{Exemplo JUnit: testando adicionar itens}

\begin{tabular}{|c|c|}
\hline
\begin{minipage}{.45\textwidth}
\begin{lstlisting}[basicstyle=\tiny]
public class Conjunto {
   private int tamanho = 0;
   
   public int tamanho() {
      return tamanho;
   }
   public void adicionar(Object o) {
      tamanho++;
   }
}
\end{lstlisting}
\end{minipage} &
\begin{minipage}{.5\textwidth}
\begin{lstlisting}[basicstyle=\tiny]
public class ConjuntoTest {
   @Test
   public void testConjuntoVazio() {
      Conjunto c = new Conjunto();
      assertEquals(0, c.tamanho());
   }

   @Test
   public void testAdicionarUm() {
      Conjunto c = new Conjunto();
      c.adicionar(new Object());
      assertEquals(1,c.tamanho());
   }
}
\end{lstlisting}
\end{minipage} \\ \hline
\end{tabular}

\vspace{2mm}

\textcolor[rgb]{0,.5,0}{\bf Passou} no teste para adicionar um elemento.

\end{frame}



\begin{frame}[fragile]{Exemplo JUnit: testando adicionar item já existente}

\begin{framed}
\begin{multicols}{2}
\begin{lstlisting}[basicstyle=\tiny]
public class ConjuntoTest {
   @Test
   public void testConjuntoVazio() {
      Conjunto c = new Conjunto();
      assertEquals(0, c.tamanho());
   }

   @Test
   public void testAdicionarUm() {
      Conjunto c = new Conjunto();
      c.adicionar(new Object());
      assertEquals(1, c.tamanho());
   }

@Test
public void testAdicionarJaExistente() {
   Conjunto c = new Conjunto();
   Object o = new Object();
   c.adicionar(o);
   c.adicionar(o);
   assertEquals(1, c.tamanho());
}

}
\end{lstlisting}
\end{multicols}
\end{framed}
\vspace{2mm}

\textcolor[rgb]{.5,0,0}{\bf Falhou} como esperado.

\end{frame}




\begin{frame}[fragile]{Exemplo JUnit: testando adicionar item já existente}

\begin{framed}
\begin{multicols}{2}
\begin{lstlisting}[basicstyle=\tiny]
public class Conjunto {
   private int tamanho = 0;
   public static final int MAX = 10;
   private Object[] items = new Object[MAX];
   
   public int tamanho() {
      return tamanho;
   }
   
   public void adicionar(Object o) {
      for (int i = 0; i < MAX; i++) {
         if (items[i] == o) {
            return ;
         }
      }
      items[tamanho++] = o;
   }
}
\end{lstlisting}
\end{multicols}
\end{framed}
\vspace{2mm}

\textcolor[rgb]{0,.5,0}{\bf Passou} como esperado. Agora que todos os testes passaram, podemos refatorar o loop.

\end{frame}




\begin{frame}[fragile]{Exemplo JUnit: refatorando}

\begin{tabular}{|c|c|}
\hline
\texttt{Conjunto} antes & \texttt{Conjunto} depois \\ \hline
\begin{minipage}{.45\textwidth}
\begin{lstlisting}[basicstyle=\tiny]
public void adicionar(Object o) {
   for (int i = 0; i < MAX; i++) {
      if (items[i] == o) {
         return ;
      }
   }
   items[tamanho++] = o;
}
\end{lstlisting}
\end{minipage} &
\begin{minipage}{.5\textwidth}
\begin{lstlisting}[basicstyle=\tiny]
public boolean contem(Object o) {
   for (int i = 0; i < MAX; i++) {
      if (items[i] == o) {
         return true;
      }
   }
   return false;
}

public void adicionar(Object o) {
   if (!contem(o)) {
      items[tamanho++] = o;
   }
}
\end{lstlisting}
\end{minipage} \\ \hline
\end{tabular}

\vspace{2mm}

\textcolor[rgb]{0,.5,0}{\bf Passou} no teste. A refatoração não quebrou o código.

\end{frame}



\begin{frame}{Exemplo JUnit: finalizando}
\begin{itemize}
\item Ainda não foi testado o limite da lista, mas isso será feito quando for feito tratamento de exceções.
\item O ciclo: testar primeiro $\longrightarrow$ desenvolver código $\longrightarrow$ verificar que testes passaram $\longrightarrow$ refatorar produziu um código com garantia de qualidade.
\item Testes podem ser facilmente executados toda vez que a classe for modificada.
\item O código estará correto se os testes tiverem sido bem definidos.
\end{itemize}
\end{frame}

\begin{frame}[fragile]{Métodos de asserção da JUnit}

\begin{center}
\begin{tabular}{|l|p{5cm}|}
\hline
\lstinline|assertTrue(teste)| & \footnotesize falha se teste booleano é \lstinline|false|.\\ \hline
\lstinline|assertFalse(teste)| & \footnotesize falha se teste booleano é \lstinline|true|.\\ \hline
\lstinline|assertEquals(esperado, teste real)| & \footnotesize falha se valores não são iguais.\\ \hline
\lstinline|assertSame(esperado, teste real)| & \footnotesize falha se valores não são os mesmos (através de \lstinline|==|).\\ \hline
\lstinline|assertNotSame(esperado, teste real)| & \footnotesize falha se valores são os mesmos (através de \lstinline|==|).\\ \hline
\lstinline|assertNull(teste)| & \footnotesize falha se valor não for \lstinline|null|.\\ \hline
\lstinline|assertNotNull(teste)| & \footnotesize falha se valor for \lstinline|null|.\\ \hline
\lstinline|fail()| & \footnotesize faz com que o teste interrompa sua execução e falhe.\\ \hline
\end{tabular}
\end{center}


\end{frame}



\begin{frame}{Melhorando a legibilidade dos testes}
\begin{itemize}
\item Anotação \lstinline|@DisplayName| define um nome para a classe ou o método de teste.
\begin{itemize}
\item Por exemplo: \lstinline|@DisplayName("Testando o metodo contem.")|
\end{itemize}

\item \lstinline|String| de informação é um parâmetro opcional que imprime informações no log de eventos durante a execução de um teste.
\begin{itemize}
  \item Será exibido caso o teste falhe.
\end{itemize}

\item Exemplo:

\begin{itemize}
\item \lstinline|assertTrue(shulambs(a. b), "chamada shulambs(" + a + "," + b + ")");|
\end{itemize}
\end{itemize}
\end{frame}



\begin{frame}[fragile]{Exemplo JUnit}

\begin{lstlisting}[basicstyle=\scriptsize]
public class TestLista {
 @Test
 public void testAddGet() {
   Lista lst = new Lista();
   lst.add(1);
   assertEquals(1, lst.get(0), "recupera 1o elemento");
   lst.add(3);
   assertEquals(3, lst.get(1), "recupera 2o elemento");
 }
 @Test
 public void testIsEmpty() {
   Lista lst = new Lista();
   assertTrue(lst.isEmpty(), "lista vazia ao criar");
   lst.add(1);
   assertFalse(lst.isEmpty(), "lista nao vazia apos inserir");
   lst.remove(0);
   assertTrue(lst.isEmpty(), "lista vazia apos remover");
 }
}
\end{lstlisting}
\end{frame}


\begin{frame}[fragile]{Inicialização e finalização}

\begin{itemize}
\item Métodos executados antes e depois de cada caso de teste. \vspace{-4mm}

\begin{lstlisting}
   @BeforeEach
   void setUp() throws Exception {
   }
   @AfterEach
   void tearDown() throws Exception {
   }
\end{lstlisting}

\item Métodos executados uma única vez, antes ou depois da execução de toda a classe de testes. \vspace{-4mm}

\begin{lstlisting}
   @BeforeAll
   static void setUpBeforeClass() throws Exception {
   }
   @AfterAll
   static void tearDownAfterClass() throws Exception {
   }
\end{lstlisting}
\end{itemize}

\end{frame}

\begin{frame}[fragile]{Exemplo JUnit: pré-configuração do teste}

\setlength{\FrameSep}{1pt}
\begin{framed}
\setlength{\columnsep}{0pt}
\begin{multicols}{2}
\begin{lstlisting}[basicstyle=\tiny]
public class ConjuntoTest {
  public static Conjunto c;

  @BeforeEach
  public void setUp() throws Exception {
    c = new Conjunto();
  }

  @Test
  public void testConjuntoVazio() {
    assertEquals(0, c.tamanho());
  }

  @Test
  public void testAdicionarUm() {
    int t = c.tamanho();
    c.adicionar(new Object());
    assertEquals(t+1, c.tamanho());
  }

  @Test
  public void testAdicionarJaExistente() {
    int t = c.tamanho();
    Object o = new Object();
    c.adicionar(o);
    c.adicionar(o);
    assertEquals(t+1, c.tamanho(), 
                 "adic. Object 2 vezes");
  }
}
\end{lstlisting}
\end{multicols}
\end{framed}

\end{frame}





\subsection{Exemplo de JUnit: Data}


\begin{frame}[fragile]{Exemplo de JUnit: Data}

{\tiny Baseado nos materiais de M. Ernst, S. Reges, D. Notkin, R. Mercer, http://www.cs.washington.edu/331/}

\begin{itemize}
\item Considere uma classe \lstinline|Data| com os seguintes métodos:
\begin{lstlisting}
public Data(int ano, int mes, int dia)
public Data()                      // hoje
public int getDia(), getMes(), getAno()
public void adicionaDias(int dias) // avanca a data em dias
public int diasNoMes()
public String diaDaSemana()        // ex. "Segunda-feira"
public boolean equals(Object o)
public boolean eAnoBisexto()
public void proximoDia()           // avanca um dia
public String toString()           // ex. 5 de marco de 2035
\end{lstlisting}

\item Proponha testes unitários para as situações a seguir:
\begin{itemize}
\item Data não pode entrar em estado inválido.
\item Função \lstinline|adicionaDias| funciona adequadamente.
\end{itemize}
\end{itemize}

\end{frame}


\begin{frame}[fragile]{Exemplo de JUnit: Data}

\begin{itemize}
\item O que está errado com o exemplo a seguir?
\end{itemize}
\vspace{-2mm}
\begin{lstlisting}
public class DataTest {
  @Test
  public void test1() {
    Data d = new Data(2050, 2, 15);
    d.adicionaDias(4);
    assertEquals(d.getAno(), 2050);
    assertEquals(d.getMes(), 2);
    assertEquals(d.getDia(), 19);
  }
  @Test
  public void test2() {
    Data d = new Data(2050, 2, 15);
    d.adicionaDias(14);
    assertEquals(d.getAno(), 2050);
    assertEquals(d.getMes(), 3);
    assertEquals(d.getDia(), 1);
  }
}
\end{lstlisting}

\end{frame}


\begin{frame}[fragile]{Exemplo de JUnit: Data}
\vspace{-2mm}
\begin{lstlisting}
public class DataTest {
  @Test
  public void test1() { // Usar nomes significativos.
    Data d = new Data(2050, 2, 15);
    d.adicionaDias(4);
    assertEquals(2050, d.getAno());  // Valor esperado
    assertEquals(2, d.getMes());     // a esquerda.
    assertEquals(19, d.getDia());    
  } // Valores esperado e real ja sao exibidos 
    // e nao precisam estar na mensagem.
  @Test
  public void test2() {
    Data d = new Data(2050, 2, 15);
    d.adicionaDias(14);
    assertEquals(2050, d.getAno(), "ano apos +14 dias");
    assertEquals(3, d.getMes(), "mes apos +14 dias");
    assertEquals(1, d.getDia(), "dia apos +14 dias");
  } // Mensagens devem explicar o que esta sendo testado
}
\end{lstlisting}

\end{frame}


\begin{frame}[fragile]{Eliminando redundâncias e tornando os testes curtos}

\begin{itemize}
\item Use objetos com valores esperados para reduzir testes.
\end{itemize} \vspace{-2mm}

\begin{lstlisting}
public class DataTest {
  @Test
  public void testAdicionaDias_numMesmoMes() {
    Data real = new Data(2050, 2, 15);
    real.adicionaDias(4);
    Data esperado = new Data(2050, 2, 19);
    assertEquals(esperado, real);  
  } // Data deve possuir metodo toString e equals.
  
  @Test                           
  public void testAdicionaDias_proximoMes() {
    Data real = new Data(2050, 2, 15);
    real.adicionaDias(14);
    Data esperado = new Data(2050, 3, 1);
    assertEquals(esperado, real, "data apos +14 dias");
  }
}
\end{lstlisting}

\end{frame}

\begin{frame}[fragile]{Eliminando redundâncias e tornando os testes curtos}
\vspace{-2mm}
\begin{lstlisting}
public class DataTest {
  @Test
  public void testAdicionaDias_numMesmoMes() {
    adicionaHelper(2050, 2, 15, 4, 2050, 2, 19);
  } 
  
  @Test
  public void testAdicionaDias_proximoMes() {
    adicionaHelper(2050, 2, 15, 14, 2050, 3, 1);
  }
  
  private void adicionaHelper(int ano1, int mes1, int dia1,
                    int adic, int ano2, int mes2, int dia2) {
    Data real = new Data(ano1, mes1, dia1);
    real.adicionaDias(adic);
    Data esper = new Data(ano2, mes2, dia2);
    assertEquals(esper, real, "data apos +" + adic + " dias");
  }
}
\end{lstlisting}

\end{frame}


\end{document}
